\chapter{\MakeUppercase{Project Implementation}}

\section{Implementation Overview}
\paragraph{}

This chapter describes how the proposed fairness-aware credit scoring framework was implemented in practice. The implementation follows a modular Python-based design so that each stage of the pipeline can be developed, tested, and extended independently.

The complete workflow is executed from a single entry script and covers:
\begin{itemize}
    \item Dataset loading
    \item Preprocessing
    \item Model creation
    \item Fairness-aware fitting
    \item Evaluation and statistical output generation
    \item Visualization
    \item Recommendation support
\end{itemize}

\section{Technology Stack and Environment}
\paragraph{}

The implementation is developed in Python and uses open-source libraries:

\begin{itemize}
    \item pandas, numpy for data handling and numerical operations
    \item scikit-learn for preprocessing, model training, and cross-validation
    \item fairlearn for fairness constraint application and fairness metrics
    \item interpret for Explainable Boosting Machine (EBM)
    \item xgboost for gradient boosted tree modeling
    \item scipy for paired t-test calculations
    \item matplotlib for plotting and chart export
\end{itemize}

Project structure is organized into reusable modules under \texttt{src}, with outputs persisted under \texttt{outputs}.

\section{Repository Structure and Module Responsibilities}
\paragraph{}

The implementation is divided into clearly scoped modules.

\subsection{Configuration Module (config.py)}
\paragraph{}

This module stores all global experiment settings:
\begin{itemize}
    \item Random seed
    \item Dataset paths and metadata
    \item Fairness levels
    \item Train/test split ratio
    \item Cross-validation folds
    \item Output paths
    \item Recommendation thresholds
\end{itemize}

Centralizing these parameters ensures consistent behavior and easy reconfiguration without modifying core logic.

\subsection{Data Loading Module (data\_loader.py)}
\paragraph{}

This module implements dataset-specific parsers and unified dataset packaging.

Key implementation points:
\begin{itemize}
    \item DatasetSpec dataclass holds dataset name, target column, protected attribute, and DataFrame
    \item \_load\_german() parses German Credit and maps encoded personal-status values to sex
    \item \_load\_australian() parses Australian dataset and ensures binary target typing
    \item \_load\_gmsc() loads CSV and creates age\_group from age threshold logic
    \item load\_datasets() reads datasets defined in config and skips missing files safely
\end{itemize}

\subsection{Preprocessing Module (preprocess.py)}
\paragraph{}

This module standardizes feature engineering for all datasets.

Implemented steps:
\begin{itemize}
    \item Split into X (features), y (target), and sensitive (protected attribute)
    \item Detect categorical and numerical columns dynamically
    \item Numeric pipeline: median imputation + standard scaling
    \item Categorical pipeline: most-frequent imputation + one-hot encoding
    \item Combine transformations using ColumnTransformer
\end{itemize}

\subsection{Model Factory Module (models.py)}
\paragraph{}

This module returns all model objects in a dictionary for iterative training.

Implemented models:
\begin{itemize}
    \item Logistic Regression
    \item Decision Tree
    \item Random Forest
    \item Gradient Boosting
    \item Explainable Boosting Machine (EBM)
    \item XGBoost
\end{itemize}

Model initialization uses a common random seed where supported.

\subsection{Fairness Module (fairness.py)}
\paragraph{}

This module contains fairness-specific implementation.

Core functions:
\begin{itemize}
    \item Fairness-level to epsilon mapping (fairness\_eps\_from\_level)
    \item Constraint-aware estimator building (build\_estimator)
    \item Conditional fitting logic for baseline vs constrained training (fit\_estimator)
    \item Probability/score extraction abstraction (predict\_scores)
    \item Fairness metric calculation (compute\_fairness\_metrics)
    \item Subgroup metric generation using MetricFrame (subgroup\_metrics)
\end{itemize}

When fairness level is 0.0, the base estimator is used directly. For levels above 0.0, the estimator is wrapped with Fairlearn Exponentiated Gradient under Demographic Parity.

\subsection{Evaluation Module (evaluation.py)}
\paragraph{}

This module implements predictive metric computation and fold summarization.

Implemented features:
\begin{itemize}
    \item Test metrics: Accuracy, F1, ROC-AUC
    \item K-fold metric extraction through compute\_cv\_metrics
    \item Fold aggregation into mean, standard deviation, and confidence interval
    \item Confidence interval estimation via normal approximation
\end{itemize}

\subsection{Visualization Module (visualize.py)}
\paragraph{}

This module generates publication-ready plots from final results tables.

Implemented plot generators:
\begin{itemize}
    \item Trade-off curve per model and dataset
    \item Pareto frontier per dataset
    \item Combined fairness-accuracy scatter across all models
    \item Baseline vs mitigated comparisons (accuracy and fairness)
\end{itemize}

\subsection{Decision Module (decision.py)}
\paragraph{}

This module implements practical configuration recommendation logic.

Selection strategy:
\begin{itemize}
    \item Try configurations satisfying maximum allowed accuracy drop threshold
    \item If configured, check explicit fairness target threshold
    \item Otherwise, choose closest fairness-performance compromise
\end{itemize}

It also generates stakeholder-readable summary text for final reporting.

\subsection{Experiment Orchestrator (main.py)}
\paragraph{}

\texttt{run\_experiment()} is the central execution routine.

High-level implemented flow:
\begin{itemize}
    \item Initialize seed and output directories
    \item Load available datasets
    \item For each dataset:
    \begin{itemize}
        \item Split train/test data
        \item Build preprocessor
        \item Iterate over all models
        \item Iterate over all fairness levels
        \item Train, evaluate, and store metrics and subgroup values
    \end{itemize}
    \item Run paired t-tests against baseline fold accuracies
    \item Save output CSV files
    \item Generate plots
    \item Print dataset-wise recommendation summaries
\end{itemize}

\section{End-to-End Execution Workflow}
\paragraph{}

The complete implementation is executed through a single script.

During execution, the system automatically:
\begin{itemize}
    \item Discovers configured datasets that exist on disk
    \item Trains all model-fairness combinations
    \item Computes predictive and fairness metrics
    \item Performs statistical testing
    \item Exports CSV artifacts and PNG figures
\end{itemize}

\section{Output Artifacts Produced by Implementation}
\paragraph{}

The implementation generates the following core artifacts:

\begin{itemize}
    \item \textbf{results.csv} -- Contains run-level metrics for each dataset-model-fairness combination
    \item \textbf{subgroup\_metrics.csv} -- Contains group-wise selection rate, TPR, and FPR values
    \item \textbf{paired\_tests.csv} -- Contains paired t-test results comparing baseline and fairness-constrained accuracy
    \item \textbf{outputs/plots/*.png and outputs/fairness\_plots\_clean/*.png} -- Visual comparative charts used for analysis and reporting
    \item \textbf{figure\_ci\_auc\_95.png} -- Confidence interval figure generated for reporting
    \item \textbf{figure\_accuracy\_heatmap.png} -- Accuracy heatmap for model vs fairness-level comparison
\end{itemize}

\section{Implementation Decisions and Rationale}
\paragraph{}

Key design decisions made in implementation include:

\begin{itemize}
    \item \textbf{Modular decomposition} -- Separate modules reduce coupling and simplify maintenance
    \item \textbf{Configuration-driven behavior} -- Experiment settings are controlled from one source (config.py), improving reproducibility
    \item \textbf{Uniform preprocessing} -- A shared preprocessing pipeline keeps model comparisons fair
    \item \textbf{Progressive fairness control} -- Multi-level fairness severity allows controlled trade-off exploration
    \item \textbf{Statistical output support} -- Fold-wise metrics and paired tests provide stronger evidence than single test scores
    \item \textbf{Artifact-first design} -- Every experiment run produces reusable tables and plots for downstream chapters
\end{itemize}

\section{Implementation Constraints and Current Scope}
\paragraph{}

The current implementation scope includes:
\begin{itemize}
    \item Group fairness under Demographic Parity constraint
    \item Tabular datasets configured in the project
    \item CPU-friendly experimentation
\end{itemize}

The current implementation does not yet include:
\begin{itemize}
    \item SHAP-based explanation pipeline integration in main run loop
    \item Calibration-difference metric computation
    \item Individual/counterfactual fairness module
\end{itemize}

These can be considered extensions without changing the overall architecture.

\section{Chapter Summary}
\paragraph{}

This chapter presented the practical realization of the proposed system as a modular, configuration-driven Python framework. The implementation covers the full experimental lifecycle from data ingestion to fairness-aware model training, evaluation, artifact generation, and recommendation output. The produced artifacts and outputs provide the foundation for the next chapter, where empirical results are analyzed and discussed.